% AUTO-GENERATED by paper/gen_tables.py -- do not edit by hand; edit the data and re-run.
\begin{table}[t]
\centering
\caption{Attack outcome per design (majority vote over all locations $\times$ trials; ties go to the worse outcome). \oS{}~succeeded, \oP{}~partial, \oB{}~blocked; for Q1, \oB{} means the node was detected and quarantined. \textit{(expected, not yet measured)}}
\label{tab:attacks}
\footnotesize
\setlength{\tabcolsep}{4pt}
\begin{tabular}{@{}llcccc@{}}
\toprule
ID & Attack & D0 & D1 & D2 & D3 \\
\midrule
A1 & Recon ping sweep & \oS & \oS & \oP & \oP \\
A2 & EHR access (HTTP) & \oS & \oS & \oS & \oB \\
A3 & IoT-server / cross-class & \oS & \oS & \oS & \oP \\
A4 & Telnet/SSH to infra & \oS & \oS & \oB & \oB \\
A5 & Rogue DHCP server & \oS & \oS & \oS & \oB \\
A6 & Unauthorised device & \oS & \oS & \oS & \oB \\
A7 & Guest $\rightarrow$ EHR pivot & \oS & \oS & \oB & \oB \\
\midrule
A8 & Rogue switch (DTP) & \oS & \oB & \oB & \oB \\
Q1 & Auto-quarantine & \oS & \oS & \oB & \oB \\
\bottomrule
\end{tabular}
\end{table}
